\documentclass[10pt]{article}
\usepackage[a4paper,margin=0.84in]{geometry}
\renewcommand{\rmdefault}{cmr}
\usepackage{amsmath,amssymb,amsthm,mathtools,amscd,booktabs,enumitem,xcolor,hyperref}
\hypersetup{colorlinks=true,linkcolor=blue!55!black,citecolor=blue!55!black,urlcolor=blue!55!black}
\setlength{\emergencystretch}{3em}
\setlength{\parindent}{0pt}
\setlength{\parskip}{0.52em}
\setlist{leftmargin=1.6em,itemsep=0.22em,topsep=0.22em}
\sloppy

\newtheorem{theorem}{Theorem}[section]
\newtheorem{proposition}[theorem]{Proposition}
\newtheorem{corollary}[theorem]{Corollary}
\newtheorem{lemma}[theorem]{Lemma}
\theoremstyle{definition}
\newtheorem{definition}[theorem]{Definition}
\newtheorem{convention}[theorem]{Convention}
\newtheorem{remark}[theorem]{Remark}

\newcommand{\Fq}{\mathbf F_q}
\newcommand{\Fbar}{\overline{\mathbf F}_q}
\newcommand{\Zl}{\mathbf Z_\ell}
\newcommand{\Ql}{\mathbf Q_\ell}
\newcommand{\Qlb}{\overline{\mathbf Q}_\ell}
\newcommand{\OO}{\mathcal O}
\newcommand{\Db}{\mathcal D}
\newcommand{\Dbc}{\mathcal D^b_c}
\newcommand{\Dctf}{\mathcal D^b_{\mathrm{ctf}}}
\newcommand{\DbcW}{\mathcal D^b_{c,\mathrm W}}
\newcommand{\DctfW}{\mathcal D^b_{\mathrm{ctf},\mathrm W}}
\newcommand{\Perf}{\operatorname{Perf}}
\newcommand{\Spec}{\operatorname{Spec}}
\newcommand{\Fix}{\operatorname{Fix}}
\newcommand{\Hom}{\operatorname{Hom}}
\newcommand{\End}{\operatorname{End}}
\newcommand{\Tr}{\operatorname{Tr}}
\newcommand{\RHom}{R\mathcal Hom}
\newcommand{\RGamma}{R\Gamma}
\newcommand{\RGammac}{R\Gamma_c}
\newcommand{\Kar}{\operatorname{Kar}}
\newcommand{\one}{\mathbf 1}
\newcommand{\depends}[1]{\par\smallskip\noindent\textit{Direct dependencies.} #1\par}

\begin{document}
\begin{center}
{\large\bfseries The Lefschetz--Verdier and Grothendieck--Lefschetz Trace Formulas\par}
\medskip
{\normalsize Fixed-point correspondences, local terms, and coefficient descent\par}
\end{center}

\begin{abstract}
The Lefschetz--Verdier formula is formulated for cohomological
correspondences with proper fixed scheme.  Its constituent statements cover
extension by zero, localization, deformation to the normal cone, purity,
fixed-point transversality, and adic reconstruction.  Scheme Frobenius and
Galois Frobenius are distinguished throughout.  Every derived category below
is an ordinary Verdier-localized triangulated category.
\end{abstract}

\tableofcontents

\section{Finite coefficients and trace finiteness}

\subsection{Geometric setting}

Fix a finite field \(k=\Fq\), an algebraic closure \(\bar k=\Fbar\), and a
prime \(\ell\ne\operatorname{char}(k)\).  Let
\[
 a_0:X_0\longrightarrow s_0:=\Spec k
\]
be separated and of finite type, and put
\[
 s:=\Spec\bar k,\qquad X:=X_0\times_{s_0}s,\qquad a:X\to s.
\]
Let \(A\) be a finite commutative \(\Zl\)-algebra.  Thus \(A\) is killed by
a power of \(\ell\).

\begin{definition}[Ambient finite-coefficient derived category]
For a scheme \(Y\) on which \(\ell\) is invertible, let
\(\mathsf{Sh}(Y,A)\) be the abelian category of sheaves of \(A\)-modules on
\(Y_{\mathrm{pro\acute et}}\).  Define
\[
 D(Y,A):=
 \mathsf{Cpx}\bigl(\mathsf{Sh}(Y,A)\bigr)
 [\{\text{quasi-isomorphisms}\}^{-1}].
\]
This is the ordinary unbounded derived category.  If \(f:Y\to Z\), derived
inverse image and derived direct image have the types
\[
 Lf^*:D(Z,A)\longrightarrow D(Y,A),\qquad
 Rf_*:D(Y,A)\longrightarrow D(Z,A),
\]
and \(Lf^*\dashv Rf_*\).
\end{definition}

\begin{definition}[Bounded constructibility and finite Tor-amplitude]
An \(A\)-module sheaf \(\mathcal F\) is constructible if \(Y\) has a finite
stratification by locally closed subschemes on which \(\mathcal F\) is
locally constant with finite \(A\)-module geometric stalks.  A complex
\(K\in D(Y,A)\) is bounded constructible if only finitely many
\(\mathcal H^i(K)\) are nonzero and each is constructible.

It has finite Tor-amplitude over \(A\) if there are integers \(r\leq s\)
such that
\[
 \mathcal H^i(K\otimes_A^L M)=0\qquad(i\notin[r,s])
\]
for every \(A\)-module sheaf \(M\) placed in degree zero.  Write
\[
 \Dctf(Y,A)\subset D(Y,A)
\]
for the full triangulated subcategory of bounded constructible complexes of
finite Tor-amplitude.
\end{definition}

\begin{definition}[Finite-coefficient Tate twists]
Choose an integer \(N\geq1\) such that \(\ell^NA=0\).  On a scheme
\(Y\) on which \(\ell\) is invertible, let \(\mu_{\ell^N,Y}\) be the
pro-etale sheaf
\[
 U\longmapsto
 \{u\in\Gamma(U,\mathcal O_U^\times):u^{\ell^N}=1\}.
\]
The rank-one invertible \(A_Y\)-module
\[
 A_Y(1):=
 \mu_{\ell^N,Y}\otimes_{\mathbf Z/\ell^N\mathbf Z}A_Y
\]
is independent, up to its canonical transition isomorphism, of replacing
\(N\) by a larger exponent annihilating \(A\).  Put
\[
 A_Y(r)=A_Y(1)^{\otimes_A r}\quad(r\geq0),
 \qquad
 A_Y(-r)=\mathcal Hom_A(A_Y(r),A_Y),
\]
and, for \(K\in D(Y,A)\), define
\[
 K(r):=K\otimes_A^LA_Y(r)\qquad(r\in\mathbf Z).
\]
Pullback carries \(A_Y(r)\) canonically to the corresponding twist on the
source, and every twist preserves bounded constructibility, perfection, and
finite Tor-amplitude.
\end{definition}

\begin{definition}[Lower shriek]
For a separated finite-type morphism \(f:Y\to Z\), choose a compactification
\[
 Y\xrightarrow{j}\overline Y\xrightarrow{\bar f}Z
\]
with \(j\) open and \(\bar f\) proper, and define
\[
 Rf_!K:=R\bar f_*j_!K.
\]
Common refinements give canonical comparison isomorphisms between choices;
they satisfy the cocycle condition and the composition isomorphism
\[
 Rg_!\,Rf_!\xrightarrow{\sim}R(gf)_!.
\]
If \(f\) is proper, the canonical map \(Rf_!\to Rf_*\) is an isomorphism.
\end{definition}

\begin{definition}[Exceptional-adjoint admissibility]
A separated finite-type morphism \(f:Y\to Z\) is
\emph{\(A\)-exceptionally admissible} if the already constructed functor
\(Rf_!:D(Y,A)\to D(Z,A)\) has a specified right adjoint
\[
 Rf^!:D(Z,A)\longrightarrow D(Y,A)
\]
together with unit and counit
\[
 \eta_f:\operatorname{id}\to Rf^!Rf_!,
 \qquad
 \varepsilon_f:Rf_!Rf^!\to\operatorname{id}.
\]
The admissibility datum also requires that \(Rf^!\) carry \(\Dctf(Z,A)\)
to \(\Dctf(Y,A)\), and supplies the composition and base-change exchange
isomorphisms involving \(Rf^!\) used below.  It is stable under the
compositions and cartesian base changes that occur in the trace diagrams.

This is an explicit hypothesis: it is available once the required finite
cohomological-dimension, amplitude, finite Tor-amplitude, and adjoint
existence theorems have been established.  No assertion that every
separated finite-type morphism is admissible is made here.
\end{definition}

\begin{theorem}[Preservation and perfectness required by traces]
Let \(f:Y\to Z\) be separated and of finite type over an algebraically
closed field of characteristic different from \(\ell\).  Then
\[
 Rf_!:\Dctf(Y,A)\longrightarrow\Dctf(Z,A)
\]
is an exact triangulated functor.  If \(f\) is proper, the same assertion
holds for \(Rf_*\).  For the structural map \(b_Y:Y\to\Spec\Omega\),
\[
 \RGammac(Y,K):=R\Gamma\bigl(\Spec\Omega,Rb_{Y!}K\bigr)
\]
belongs to \(\Perf(A)\) for every \(K\in\Dctf(Y,A)\).  Every derived
geometric stalk \(K_{\bar y}\) also belongs to \(\Perf(A)\).
\depends{Constructibility of finite-torsion direct image, finite
cohomological dimension, preservation of finite Tor-amplitude, and the
compactification definition of \(Rf_!\).}
\end{theorem}

\section{The two Frobenius operations}

\begin{definition}[Galois Frobenius]
Let
\[
 \sigma:\bar k\to\bar k,\quad z\mapsto z^q,
 \qquad \gamma:=\sigma^{-1}.
\]
The automorphism \(\sigma\) is arithmetic Galois Frobenius and \(\gamma\)
is geometric Galois Frobenius.  Base change in the \(\bar k\)-factor gives
semilinear inverse automorphisms
\[
 g_\sigma,g_\gamma:X\longrightarrow X.
\]
On \(A_0\otimes_k\bar k\), for \(X_0=\Spec A_0\), they act as
\[
 a\otimes b\mapsto a\otimes b^q,
 \qquad
 a\otimes b\mapsto a\otimes\gamma(b).
\]
\end{definition}

\begin{definition}[Scheme Frobenius]
The \(q\)-power scheme Frobenius is the \(k\)-morphism
\(\Phi_0:X_0\to X_0\).  Its base change
\[
 \Phi:=\Phi_0\times_{s_0}\operatorname{id}_s:X\to X
\]
is a \(\bar k\)-morphism and on \(A_0\otimes_k\bar k\) is
\[
 \Phi^\#(a\otimes b)=a^q\otimes b.
\]
It is finite and radicial, hence proper and a universal homeomorphism.  It
is generally not an automorphism.
\end{definition}

\begin{proposition}[Comparison of scheme and Galois pullback]
Absolute \(q\)-power Frobenius on \(X\) factors as
\[
 \operatorname{Fr}^{(q)}_X=g_\sigma\Phi=\Phi g_\sigma.
\]
Topological invariance for universal homeomorphisms gives a natural
isomorphism
\[
 \kappa:L\Phi^*\xrightarrow{\sim}Lg_\gamma^*
\]
on \(D(X,A)\).  For every \(m\geq1\) there is a multiplicative iterate
\[
 \kappa^{(m)}:L(\Phi^m)^*\xrightarrow{\sim}L(g_\gamma^m)^*.
\]
\depends{The displayed factorization and topological invariance of the
pro-etale topos under universal homeomorphisms.}
\end{proposition}

\section{Strict Weil complexes and Frobenius operators}

\begin{definition}[Strict Weil category]
Let \(\mathsf{Sh}_{\mathrm W}(X,A)\) be the abelian category of pairs
\[
 (\mathcal F,\varphi_{\mathcal F}),\qquad
 \varphi_{\mathcal F}:g_\gamma^*\mathcal F\xrightarrow{\sim}\mathcal F,
\]
with morphisms commuting with \(\varphi\).  A strict Weil complex is a
cochain complex in this abelian category.  Localizing at equivariant
quasi-isomorphisms gives
\[
 D_{\mathrm W}(X,A):=
 D\bigl(\mathsf{Sh}_{\mathrm W}(X,A)\bigr).
\]
The full triangulated subcategory whose underlying complex lies in
\(\Dctf(X,A)\) is denoted \(\DctfW(X,A)\).
\end{definition}

\begin{definition}[Iterated Weil isomorphism]
For \((K,\varphi)\in\DctfW(X,A)\), define
\[
 \varphi^{[m]}:L(g_\gamma^m)^*K\xrightarrow{\sim}K
\]
by
\[
 \varphi^{[m]}=
 \varphi\circ Lg_\gamma^*\varphi\circ\cdots\circ
 L(g_\gamma^{m-1})^*\varphi.
\]
The corresponding scheme-Frobenius morphism is
\[
 u_m:=\varphi^{[m]}\circ\kappa_K^{(m)}:
 L(\Phi^m)^*K\longrightarrow K.
\]
The cocycle identity is
\[
 \varphi^{[m+n]}
 =\varphi^{[m]}\circ L(g_\gamma^m)^*\varphi^{[n]}
\]
after the canonical pullback identifications.
\end{definition}

\begin{definition}[Global geometric Frobenius]
Compact-support base change for the automorphism \(g_\gamma^m\) gives
\[
 g_{\gamma,c}^{m,*}:
 \RGammac(X,K)\xrightarrow{\sim}
 \RGammac\bigl(X,L(g_\gamma^m)^*K\bigr).
\]
Define
\[
 F_{c,K}^{\,m}:=
 \RGammac(X,\varphi^{[m]})\circ g_{\gamma,c}^{m,*}.
\]
It is an endomorphism of the perfect complex \(\RGammac(X,K)\).
\end{definition}

\begin{definition}[Local geometric Frobenius]
For \(y\in X_0(\mathbf F_{q^m})\), choose the associated geometric point
\(\bar y:\Spec\bar k\to X\).  The equality
\(g_\gamma^m\bar y=\bar y\Spec(\gamma^m)\) gives a canonical stalk
transport
\[
 \tau_{y,m}:K_{\bar y}\xrightarrow{\sim}
 \bigl(L(g_\gamma^m)^*K\bigr)_{\bar y}.
\]
Define
\[
 F_{y,m}:=
 (\varphi^{[m]})_{\bar y}\circ\tau_{y,m}:K_{\bar y}\to K_{\bar y}.
\]
Changing the lift \(\bar y\) conjugates this endomorphism.

If \(x\in|X_0|\) has degree \(d_x=[\kappa(x):\Fq]\), one circuit around
\(x\) defines \(F_x:K_{\bar x}\to K_{\bar x}\).  When \(d_x\mid m\), the
operator attached to a rational point above \(x\) is conjugate to
\[
 F_x^{\,m/d_x},
\]
not to \(F_x^m\).
\end{definition}

\section{Perfect complexes and categorical trace}

\begin{definition}[Perfect complex]
For a commutative ring \(R\), let \(\Perf(R)\subset D(R)\) be the full
triangulated subcategory of objects represented by bounded complexes of
finitely generated projective \(R\)-modules.  If \(P\in\Perf(R)\), put
\[
 P^\vee:=R\Hom_R(P,R).
\]
There are coevaluation and evaluation maps
\[
 \operatorname{coev}_P:R\to P\otimes_R^LP^\vee,\qquad
 \operatorname{ev}_P:P^\vee\otimes_R^LP\to R
\]
satisfying the two triangle identities.
\end{definition}

\begin{definition}[Categorical trace]
For \(v:P\to P\) with \(P\in\Perf(R)\), define
\[
 \operatorname{tr}_R(v\mid P)\in\End_{D(R)}(R)\cong R
\]
by the composite
\[
\begin{aligned}
R&\xrightarrow{\operatorname{coev}_P}P\otimes_R^LP^\vee
\xrightarrow{v\otimes1}P\otimes_R^LP^\vee\\
&\xrightarrow{\mathrm{sym}}P^\vee\otimes_R^LP
\xrightarrow{\operatorname{ev}_P}R.
\end{aligned}
\]
If \(R\) is a field, this equals
\[
 \sum_i(-1)^i\Tr\bigl(H^i(v)\mid H^i(P)\bigr).
\]
For a nonfield ring, the categorical definition is used; no projectivity of
the individual cohomology modules is required.
\end{definition}

\begin{proposition}[Formal laws for categorical trace]
Categorical trace is invariant under conjugacy.  If
\[
 P'\to P\to P''\to P'[1]
\]
is a distinguished triangle in \(\Perf(R)\) and
\((v',v,v'')\) is an endomorphism of that triangle, then
\[
 \operatorname{tr}_R(v\mid P)=
 \operatorname{tr}_R(v'\mid P')+
 \operatorname{tr}_R(v''\mid P'').
\]
For a ring homomorphism \(\rho:R\to R'\),
\[
 \rho\bigl(\operatorname{tr}_R(v\mid P)\bigr)=
 \operatorname{tr}_{R'}
 \bigl(v\otimes_R^LR'\mid P\otimes_R^LR'\bigr).
\]
\depends{Dualizability of perfect complexes, the axioms of a symmetric
monoidal triangulated category, and derived scalar extension.}
\end{proposition}

\section{Cohomological correspondences and trace classes}

In this section \(\Omega\) is algebraically closed of characteristic
different from \(\ell\), \(X\) is separated of finite type over
\(\Omega\), and \(A\) is a finite coefficient ring.  Put
\(s_\Omega=\Spec\Omega\).  Every morphism to which \(R(-)^!\) is applied
below---in particular the structural maps, exceptional correspondence legs,
and their fixed-point base changes---is assumed \(A\)-exceptionally
admissible, with the composition and exchange data of the preceding
definition.

\begin{definition}[Cohomological correspondence]
Let \(C\) be separated of finite type over \(\Omega\), and let
\[
 c=(c_1,c_2):C\longrightarrow X\times_\Omega X.
\]
Assume that \(c_1\) is \(A\)-exceptionally admissible.
For \(K\in\Dctf(X,A)\), a cohomological correspondence along \(c\) is a
morphism
\[
 u:Lc_2^*K\longrightarrow Rc_1^!K
 \quad\text{in }D(C,A).
\]
Thus \(c_2\) is the ordinary pullback leg and \(c_1\) is the exceptional
pullback leg.
\end{definition}

\begin{definition}[Global operator of a correspondence]
Assume that \(c_2\) is proper.  The global operator
\[
 T_c(u):\RGammac(X,K)\longrightarrow\RGammac(X,K)
\]
is
\[
\begin{aligned}
\RGammac(X,K)
&\xrightarrow{\eta^*_{c_2}}
\RGammac\bigl(X,Rc_{2*}Lc_2^*K\bigr)
\xrightarrow{\sim}\RGammac\bigl(C,Lc_2^*K\bigr)\\
&\xrightarrow{\RGammac(C,u)}
\RGammac\bigl(C,Rc_1^!K\bigr)
\xrightarrow{\sim}\RGammac\bigl(X,Rc_{1!}Rc_1^!K\bigr)\\
&\xrightarrow{\RGammac(X,\varepsilon_{c_1})}
\RGammac(X,K).
\end{aligned}
\]
The first arrow is induced by the unit of \(Lc_2^*\dashv Rc_{2*}\); the
first isomorphism uses \(Rc_{2*}\simeq Rc_{2!}\); the last arrow uses the
counit of \(Rc_{1!}\dashv Rc_1^!\).
\end{definition}

\begin{definition}[Fixed scheme]
Let \(\Delta_X:X\to X\times X\) be the diagonal.  Define
\[
 F:=\Fix(c):=
 C\times_{X\times X,\,c,\,\Delta_X}X.
\]
Let \(i_F:F\to C\) be the first projection and
\(d:F\to X\) the second projection.  Then
\[
 c_1i_F=d=c_2i_F.
\]
\end{definition}

\begin{definition}[Verdier duality data]
For a separated finite-type \(\Omega\)-scheme \(Y\) whose structural map
\(a_Y:Y\to s_\Omega\) is \(A\)-exceptionally admissible, put
\[
 \omega_{Y,A}:=Ra_Y^!A_{s_\Omega},\qquad
 D_YM:=R\mathcal Hom_A(M,\omega_{Y,A}).
\]
For \(M\in\Dctf(Y,A)\), Verdier biduality gives
\[
 M\xrightarrow{\sim}D_YD_YM,
\]
and evaluation is
\[
 \operatorname{ev}_M:D_YM\otimes_A^LM\to\omega_{Y,A}.
\]
\end{definition}

\begin{proposition}[Fixed-point evaluation morphism]
For \(K\in\Dctf(X,A)\), define
\[
 \mathcal H_c(K):=
 R\mathcal Hom_A(Lc_2^*K,Rc_1^!K)\in D(C,A).
\]
There is a natural morphism
\[
 \theta_{c,K}:Li_F^*\mathcal H_c(K)\longrightarrow\omega_{F,A}.
\]
It is obtained from Verdier biduality, tensor--internal-Hom adjunction, the
cartesian fixed square, the equality \(c_1i_F=c_2i_F=d\), and evaluation
\[
 D_XK\otimes_A^LK\longrightarrow\omega_{X,A},
\]
followed by \(Rd^!\omega_{X,A}\simeq\omega_{F,A}\).  It is natural in
\((C,c,K)\), compatible with shifts, and \(A\)-linear.
\depends{Verdier biduality, exceptional base change, the projection formula,
and composition for exceptional inverse image.}
\end{proposition}

\begin{definition}[Verdier trace class]
A correspondence \(u:Lc_2^*K\to Rc_1^!K\) is equivalently a section
\[
 \widetilde u:A_C\longrightarrow\mathcal H_c(K).
\]
Its Verdier trace class is
\[
 \operatorname{Tr}_c(u):
 A_F\xrightarrow{Li_F^*\widetilde u}
 Li_F^*\mathcal H_c(K)
 \xrightarrow{\theta_{c,K}}\omega_{F,A}.
\]
Thus
\[
 \operatorname{tr}_c:
 \Hom_{D(C,A)}(Lc_2^*K,Rc_1^!K)
 \longrightarrow
 \Hom_{D(F,A)}(A_F,\omega_{F,A})
\]
is a natural \(A\)-linear map.
\end{definition}

\begin{definition}[Integration and local term]
Assume \(F\) is proper over \(\Omega\).  The counit of
\(Ra_{F!}\dashv Ra_F^!\), together with \(Ra_{F!}\simeq Ra_{F*}\), defines
\[
 \int_F:
 \Hom_{D(F,A)}(A_F,\omega_{F,A})\longrightarrow A.
\]
If \(\beta\subset F\) is an open-and-closed connected component, define
\[
 \operatorname{LT}_\beta(u):=
 \int_\beta\bigl(\operatorname{Tr}_c(u)|_\beta\bigr)\in A.
\]
Equivalently, it is the endomorphism of \(A\) represented by
\[
\begin{aligned}
A&\longrightarrow R\Gamma(\beta,A_\beta)
\xrightarrow{R\Gamma(\operatorname{Tr}_c(u)|_\beta)}
R\Gamma(\beta,\omega_{\beta,A})
\xrightarrow{\varepsilon_{a_\beta}}A.
\end{aligned}
\]
\end{definition}

\section{Finite-coefficient Lefschetz--Verdier formalism}

Throughout this section and the Frobenius-graph section below, retain the
exceptional-adjoint admissibility hypotheses stated in the construction of
trace classes.  Thus every displayed use of \(R(-)^!\), Verdier duality, or
exceptional base change is made only with the corresponding specified
adjunction and exchange data.

\subsection{Functoriality and localization}

\begin{proposition}[Additivity of trace classes]
Let
\[
 K'\longrightarrow K\longrightarrow K''\longrightarrow K'[1]
\]
be a distinguished triangle in \(\Dctf(X,A)\).  Suppose
\((u',u,u'')\) is a morphism from the triangle obtained by \(Lc_2^*\) to
the triangle obtained by \(Rc_1^!\).  Then
\[
 \operatorname{Tr}_c(u)=
 \operatorname{Tr}_c(u')+\operatorname{Tr}_c(u'')
\]
in \(\Hom_{D(F,A)}(A_F,\omega_{F,A})\).  If \(F\) is proper, the same
identity holds for every local term; if \(c_2\) is proper, it also holds for
the categorical traces of the three global operators.
\depends{Linearity of \(\theta_{c,K}\), compatibility of Verdier duality
with distinguished triangles, and additivity of categorical trace.}
\end{proposition}

\begin{definition}[Stable open--closed correspondence datum]
Let \(j:U\hookrightarrow X\) be an open immersion with closed complement
\(i:Z\hookrightarrow X\).  It is \emph{stable for} \(c\) if the two open
inverse images and the two closed inverse images are identified as
subschemes of \(C\):
\[
 C_U:=c_1^{-1}(U)=c_2^{-1}(U),
 \qquad
 C_Z:=c_1^{-1}(Z)=c_2^{-1}(Z).
\]
Writing \(j_C:C_U\hookrightarrow C\) and \(i_C:C_Z\hookrightarrow C\),
define
\[
 c_U=(c_{1,U},c_{2,U}):C_U\to U\times U,
 \qquad
 c_Z=(c_{1,Z},c_{2,Z}):C_Z\to Z\times Z.
\]
This definition does not assert that \(\Fix(c_U)\) and \(\Fix(c_Z)\) are
open-and-closed components of \(\Fix(c)\).
\end{definition}

\begin{definition}[Open extension-by-zero trace datum]
Let \(j_0:U_0\hookrightarrow X_0\) be open with closed complement
\(i_0:Z_0\hookrightarrow X_0\), and let \(j:U\hookrightarrow X\) and
\(i:Z\hookrightarrow X\) be their geometric base changes.  For
\(K\in\DctfW(X,A)\), the open--closed datum consists of the localization
arrows
\[
 e_{j,K}:Rj_!Lj^*K\longrightarrow K,
 \qquad
 r_{i,K}:K\longrightarrow Ri_*Li^*K,
\]
their connecting arrow
\(Ri_*Li^*K\to Rj_!Lj^*K[1]\), and the compact-support composition
isomorphism
\[
 \chi_j:Ra_{X!}Rj_!\xrightarrow{\sim}Ra_{U!},
 \qquad a_U=a_Xj.
\]
Here \(e_{j,K}\) is the counit of \(j_!\dashv j^*\), while \(r_{i,K}\)
is the unit of \(i^*\dashv i_*\).  Their Frobenius structures are those
induced by the exchange maps stated below.
\end{definition}

\begin{lemma}[Underived extension by zero has zero complement stalks]
For an \(A\)-module sheaf \(\mathcal L\) on \(U\), extension by zero is the
left adjoint
\[
 j_!: \mathsf{Sh}(U,A)\longrightarrow\mathsf{Sh}(X,A)
\]
to \(j^*\), and its geometric stalks are
\[
 (j_!\mathcal L)_{\bar x}\cong
 \begin{cases}
  \mathcal L_{\bar x},&\bar x\text{ factors through }U,\\
  0,&\bar x\text{ factors through }Z.
 \end{cases}
\]
Thus, for a bounded complex \(L^\bullet\) of constructible finite
\(A\)-module sheaves and \(\bar z\to Z\), the stalk complex
\((j_!L^\bullet)_{\bar z}\) is degreewise zero.
\depends{The sheaf-theoretic extension-by-zero adjunction and its geometric
stalk formula.}
\end{lemma}

\begin{proposition}[Finite-derived zero complement restriction]
Extension by zero is exact, preserves constructibility and finite
Tor-amplitude, and therefore induces the exact triangulated functor
\[
 Rj_!=j_!:\Dctf(U,A)\longrightarrow\Dctf(X,A)
\]
represented by termwise extension by zero.  For every
\(L\in\Dctf(U,A)\), naturally in \(L\),
\[
 Lj^*Rj_!L\xrightarrow{\sim}L,
 \qquad
 Li^*Rj_!L\xrightarrow{\sim}0.
\]
Equivalently, every derived geometric stalk of \(Rj_!L\) on \(Z\) is the
zero object of \(D(A)\).
\depends{Exactness of extension by zero, a bounded constructible
representative, the underived zero-stalk lemma, and termwise computation of
inverse image.}
\end{proposition}

\begin{definition}[Finite Frobenius fixed-point set]
For \(T_0\in\{X_0,U_0,Z_0\}\), let \(T=T_0\times_{\Fq}\bar k\), and let
\(\Phi_T:T\to T\) be scheme Frobenius.  For \(m\geq1\), define the finite
set
\[
 \Fix_m(T_0):=
 \{\bar t\in T(\bar k):\Phi_T^m(\bar t)=\bar t\}.
\]
It is canonically identified with \(T_0(\mathbf F_{q^m})\).  This is a set
of geometric fixed points; it is not the assertion that an arbitrary
correspondence fixed scheme splits into open-and-closed pieces.
\end{definition}

\begin{lemma}[Frobenius fixed points split across the open--closed pair]
For every \(m\geq1\), the inclusions induce a canonical bijection of finite
sets
\[
 \Fix_m(X_0)\xrightarrow{\sim}
 \Fix_m(U_0)\amalg\Fix_m(Z_0).
\]
\depends{Complementarity of \(U_0\) and \(Z_0\), their stability under
scheme Frobenius, and the rational-point interpretation of Frobenius fixed
points.}
\end{lemma}

\begin{lemma}[A Frobenius local sum splits across the two supports]
For a function \(b:\Fix_m(X_0)\to A\), let \(b_U\) and \(b_Z\) be its
restrictions.  Then
\[
 \sum_{x\in\Fix_m(X_0)}b(x)=
 \sum_{x\in\Fix_m(U_0)}b_U(x)+
 \sum_{x\in\Fix_m(Z_0)}b_Z(x).
\]
The same formula holds for values in any commutative coefficient ring.
\depends{The preceding finite-set bijection and reindexing of a finite sum
over a disjoint union.}
\end{lemma}

\begin{proposition}[Underived extension by zero respects Frobenius pullback]
Let \(h\) be either scheme Frobenius \(\Phi_X\) or Galois Frobenius
\(g_\gamma\), and let \(h_U\) be its restriction to \(U\).  For every
constructible finite \(A\)-module sheaf \(\mathcal L\) on \(U\), the
cartesian open square induces a natural isomorphism
\[
 \beta^{0}_{h,\mathcal L}:h^*j_!\mathcal L
 \xrightarrow{\sim}j_!h_U^*\mathcal L.
\]
These isomorphisms commute with morphisms of sheaves and with composition
of the Frobenius maps.
\depends{Underived base change for extension by zero in the cartesian open
square and its sectionwise inverse.}
\end{proposition}

\begin{proposition}[Finite-derived extension by zero respects Frobenius]
For \(h=\Phi_X\) or \(g_\gamma\), termwise derivation of the preceding
isomorphism gives a natural isomorphism of exact triangulated functors
\[
 \beta_{h}:Lh^*Rj_!\xrightarrow{\sim}Rj_!Lh_U^*
 \quad:\quad\Dctf(U,A)\longrightarrow\Dctf(X,A).
\]
It is compatible with every positive iterate.  Hence a strict Weil
isomorphism
\(\varphi_U:L(g_{\gamma,U})^*K_U\xrightarrow{\sim}K_U\) extends as
\[
 Lg_\gamma^*Rj_!K_U
 \xrightarrow{\beta_{g_\gamma}}Rj_!L(g_{\gamma,U})^*K_U
 \xrightarrow{Rj_!\varphi_U}Rj_!K_U,
\]
and its \(m\)-th scheme-Frobenius arrow is
\[
 L(\Phi_X^m)^*Rj_!K_U
 \xrightarrow{\beta_{\Phi_X^m}}Rj_!L(\Phi_U^m)^*K_U
 \xrightarrow{Rj_!u_{m,U}}Rj_!K_U.
\]
The comparison
\[
 \RGammac(X,Rj_!K_U)\xrightarrow{\sim}\RGammac(U,K_U)
\]
induced by \(\chi_j\) intertwines these Frobenius operators.  At a point
of \(U_0(\mathbf F_{q^m})\) their stalk traces agree; at a point of
\(Z_0(\mathbf F_{q^m})\) the stalk complex and its trace are zero.
\depends{The underived Frobenius exchange, exactness of extension by zero,
the strict Weil iterate cocycle, the scheme--Galois comparison, and
composition for lower shriek.}
\end{proposition}

\begin{proposition}[Extension by zero of a cohomological correspondence]
Assume that \((U,Z)\) is stable for \(c\).  If
\[
 K_U\in\Dctf(U,A),\qquad
 u_U:Lc_{2,U}^*K_U\longrightarrow Rc_{1,U}^!K_U,
\]
define its extension by the fully typed composite
\[
\begin{aligned}
 j_!u_U:\quad Lc_2^*Rj_!K_U
 &\xrightarrow{\sim}Rj_{C!}Lc_{2,U}^*K_U
 \xrightarrow{Rj_{C!}u_U}Rj_{C!}Rc_{1,U}^!K_U\\
 &\xrightarrow{\sim}Rc_1^!Rj_!K_U.
\end{aligned}
\]
Put \(F=\Fix(c)\), \(F_U=\Fix(c_U)\), and \(F_Z=\Fix(c_Z)\), and write
\(\jmath_F:F_U\hookrightarrow F\) and \(\imath_F:F_Z\hookrightarrow F\)
for the induced immersions.  Naturality of fixed-point evaluation gives
\[
 L\jmath_F^*\operatorname{Tr}_c(j_!u_U)
   =\operatorname{Tr}_{c_U}(u_U),
 \qquad
 L\imath_F^*\operatorname{Tr}_c(j_!u_U)=0.
\]
\depends{Stable open exchange for the two legs, finite-derived zero
complement restriction, and naturality of the fixed-point evaluation map.}
\end{proposition}

\begin{proposition}[Closed extension of a cohomological correspondence]
Assume that \((U,Z)\) is stable for \(c\).  For
\[
 K_Z\in\Dctf(Z,A),\qquad
 u_Z:Lc_{2,Z}^*K_Z\longrightarrow Rc_{1,Z}^!K_Z,
\]
define its closed extension on \(Ri_*K_Z\) by the composite
\[
\begin{aligned}
 i_*u_Z:\quad Lc_2^*Ri_*K_Z
 &\xrightarrow{\sim}Ri_{C*}Lc_{2,Z}^*K_Z
 \xrightarrow{Ri_{C*}u_Z}Ri_{C*}Rc_{1,Z}^!K_Z\\
 &\xrightarrow{\sim}Rc_1^!Ri_*K_Z.
\end{aligned}
\]
The first isomorphism is proper base change for the cartesian square defined
by \(c_2\) and \(i\); the last is exceptional exchange for the cartesian
square defined by \(c_1\) and \(i\).  With
\(F=\Fix(c)\), \(F_U=\Fix(c_U)\), and \(F_Z=\Fix(c_Z)\), the induced
immersions \(\jmath_F:F_U\hookrightarrow F\) and
\(\imath_F:F_Z\hookrightarrow F\) satisfy
\[
 L\imath_F^*\operatorname{Tr}_c(i_*u_Z)
   =\operatorname{Tr}_{c_Z}(u_Z),
 \qquad
 L\jmath_F^*\operatorname{Tr}_c(i_*u_Z)=0.
\]
\depends{Proper base change for a closed immersion, exceptional exchange,
finite-derived zero restriction to the complementary open, and naturality
of the fixed-point evaluation morphism.}
\end{proposition}

\begin{proposition}[Extension by zero and open--closed localization]
For a stable open--closed correspondence datum, the triangle
\[
 Rj_!Lj^*K\xrightarrow{e_{j,K}}K
 \xrightarrow{r_{i,K}}Ri_*Li^*K
 \longrightarrow Rj_!Lj^*K[1]
\]
is distinguished in \(\Dctf(X,A)\).  For a correspondence
\(u:Lc_2^*K\to Rc_1^!K\), put \(K_U=Lj^*K\) and \(K_Z=Li^*K\).
Ordinary and exceptional exchange give restricted correspondences
\[
 u_U:Lc_{2,U}^*K_U\longrightarrow Rc_{1,U}^!K_U,
 \qquad
 u_Z:Lc_{2,Z}^*K_Z\longrightarrow Rc_{1,Z}^!K_Z.
\]
The arrows \(j_!u_U\), \(u\), and \(i_*u_Z\) form a morphism from the
\(Lc_2^*\)-image of the localization triangle to its \(Rc_1^!\)-image.
The associated trace classes satisfy
\[
 \operatorname{Tr}_c(u)=
 \operatorname{Tr}_c(j_!u_U)+\operatorname{Tr}_c(i_*u_Z)
 \quad\text{in }
 \Hom_{D(F,A)}(A_F,\omega_{F,A}).
\]
If \(F\) is proper, integration preserves this equality.  If \(c_2\) is
proper, categorical trace gives
\[
\begin{aligned}
 \operatorname{tr}_A\bigl(T_c(u)\mid\RGammac(X,K)\bigr)
 &=\operatorname{tr}_A\bigl(T_c(j_!u_U)\mid
       \RGammac(X,Rj_!Lj^*K)\bigr)\\
 &\quad+\operatorname{tr}_A\bigl(T_c(i_*u_Z)\mid
       \RGammac(X,Ri_*Li^*K)\bigr).
\end{aligned}
\]
For the graph of \(\Phi_X^m\), the local-term equality is exactly the
finite fixed-point sum split above; the extension-by-zero term vanishes on
\(Z_0(\mathbf F_{q^m})\).  Every displayed identity commutes with
Frobenius.
\depends{The recollement triangle, stable exchange for the two
correspondence legs, additivity of Verdier trace classes and categorical
trace, the finite fixed-point sum split, and derived Frobenius
compatibility of extension by zero.}
\end{proposition}

\begin{corollary}[Finite-stratification devissage]
Let
\[
 \varnothing=X_{-1}\subset X_0\subset\cdots\subset X_N=X
\]
be a finite filtration by closed subschemes stable under \(c\), and put
\(S_r=X_r\setminus X_{r-1}\).  If the trace identity holds for the
restriction of \((K,u)\) to every \(S_r\), it holds on \(X\).  If the
filtration and correspondence descend to \(\Fq\), the reduction is
compatible with every Frobenius iterate.
\depends{Extension by zero, open--closed localization, and induction on
the finite length of the filtration.}
\end{corollary}

\begin{definition}[Proper morphism of correspondences and its pushforward]
Let \(f:X\to Y\) and \(g:C\to D\) be proper, and let
\[
 c=(c_1,c_2):C\to X\times X,
 \qquad d=(d_1,d_2):D\to Y\times Y
\]
satisfy \(d_rg=fc_r\) for \(r=1,2\), with both leg squares cartesian.  For
\(K\in\Dctf(X,A)\) and
\(u:Lc_2^*K\to Rc_1^!K\), define
\[
\begin{aligned}
 f_*u:Ld_2^*Rf_*K
 &\xrightarrow{\sim}Rg_*Lc_2^*K
 \xrightarrow{Rg_*u}Rg_*Rc_1^!K
 \xrightarrow{\sim}Rd_1^!Rf_*K.
\end{aligned}
\]
The first isomorphism is proper base change and the last is the proper
duality exchange isomorphism.  The cartesian data induce a proper map
\[
 g_F:F_c:=\Fix(c)\longrightarrow F_d:=\Fix(d).
\]
\end{definition}

\begin{definition}[Proper pushforward of a trace class]
Composition of structural maps identifies
\[
 \omega_{F_c,A}\simeq Rg_F^!\omega_{F_d,A}.
\]
Since \(g_F\) is proper, define its dualizing trace by
\[
 \operatorname{Tr}_{g_F}^{\omega}:
 Rg_{F*}\omega_{F_c,A}
 \xrightarrow{\sim}Rg_{F!}Rg_F^!\omega_{F_d,A}
 \xrightarrow{\varepsilon_{g_F}}\omega_{F_d,A}.
\]
For
\(\alpha:A_{F_c}\to\omega_{F_c,A}\), its proper pushforward is the
following complete composite, including the adjunction unit:
\[
\begin{aligned}
 (g_F)_*\alpha:\quad A_{F_d}
 &\xrightarrow{\eta^*_{g_F}}
 Rg_{F*}Lg_F^*A_{F_d}
 \xrightarrow{\sim}Rg_{F*}A_{F_c}\\
 &\xrightarrow{Rg_{F*}\alpha}Rg_{F*}\omega_{F_c,A}
 \xrightarrow{\operatorname{Tr}_{g_F}^{\omega}}\omega_{F_d,A}.
\end{aligned}
\]
Thus \((g_F)_*\) has the type
\[
 \Hom_{D(F_c,A)}(A_{F_c},\omega_{F_c,A})
 \longrightarrow
 \Hom_{D(F_d,A)}(A_{F_d},\omega_{F_d,A}).
\]
\end{definition}

\begin{theorem}[Proper pushforward compatibility of Verdier trace classes]
For the proper morphism of correspondences above,
\[
 \operatorname{Tr}_d(f_*u)=(g_F)_*\operatorname{Tr}_c(u)
 \quad\text{in }\Hom_{D(F_d,A)}(A_{F_d},\omega_{F_d,A}).
\]
\depends{Proper base change, the proper duality exchange isomorphism,
projection, and compatibility of fixed-point evaluation with the
adjunction unit \(A_{F_d}\to Rg_{F*}A_{F_c}\) and the dualizing counit.}
\end{theorem}

\begin{corollary}[Integration and local terms under proper pushforward]
If \(F_d\) is proper, then \(F_c\) is proper and every trace class
\(\alpha:A_{F_c}\to\omega_{F_c,A}\) satisfies
\[
 \int_{F_d}(g_F)_*\alpha=\int_{F_c}\alpha.
\]
For an open-and-closed connected component \(\beta\subset F_d\),
\[
 \operatorname{LT}_\beta(f_*u)=
 \sum_{\substack{\alpha_0\in\pi_0(F_c)\\g_F(\alpha_0)\subset\beta}}
 \operatorname{LT}_{\alpha_0}(u).
\]
In particular, integration of the proper-pushforward trace identity gives
the same scalar in \(A\) on source and target.
\depends{The preceding trace-class compatibility, transitivity of the
exceptional counit, and additivity of integration on open-and-closed
components.}
\end{corollary}

\begin{proposition}[Etale pullback and excision]
Let \(e:X'\to X\) be etale, let \(c'\) be the pulled-back correspondence,
and let \(u'\) be the pullback of \(u\).  Then
\[
 \omega_{X',A}\simeq Le^*\omega_{X,A},
\qquad
 \operatorname{Tr}_{c'}(u')=
 Le_{\Fix}^*\operatorname{Tr}_c(u).
\]
In particular, the local term at an isolated fixed point depends only on an
arbitrarily small etale neighborhood of that point.
\depends{Etale purity \(Re^!\simeq Le^*\), base change, and naturality of
evaluation.}
\end{proposition}

\subsection{Deformation and the transverse local term}

\begin{definition}[Deformation to the normal cone]
Let \(\Delta_X:X\hookrightarrow X\times X\).  The deformation space
\[
 \operatorname{Def}_{\Delta_X}(X\times X)\longrightarrow\mathbf A^1
\]
is the open complement of the strict transform of
\((X\times X)\times\{0\}\) in the blow-up of
\((X\times X)\times\mathbf A^1\) along
\(\Delta_X\times\{0\}\).  Its fiber over \(1\) is \(X\times X\), and its
fiber over \(0\) is the normal cone
\[
 C_{\Delta_X}(X\times X).
\]
\end{definition}

\begin{definition}[Deformation correspondence and its specialization map]
Put
\[
 \widetilde X=X\times_\Omega\mathbf A^1,
 \qquad
 \mathfrak D=\operatorname{Def}_{\Delta_X}(X\times X),
 \qquad
 \pi_{\mathfrak D}:\mathfrak D\to\mathbf A^1.
\]
Functoriality of deformation to the normal cone for the projections
\(p_r:X\times X\to X\) gives morphisms over \(\mathbf A^1\)
\[
 \widetilde p_r:\mathfrak D\longrightarrow\widetilde X
 \qquad(r=1,2).
\]
A deformation datum for
\[
 c=(c_1,c_2):C\longrightarrow X\times X,
 \qquad
 u:Lc_2^*K\longrightarrow Rc_1^!K
\]
consists of the following typed objects:
\begin{enumerate}[label=\textup{(\alph*)}]
\item a separated finite-type morphism
\(\rho:\widetilde C\to\mathbf A^1\) and a morphism over
\(\mathbf A^1\)
\[
 d:\widetilde C\longrightarrow\mathfrak D;
\]
\item the two legs
\[
 \widetilde c_r:=\widetilde p_r d:\widetilde C\longrightarrow\widetilde X,
 \qquad
 \widetilde c=(\widetilde c_1,\widetilde c_2):
 \widetilde C\longrightarrow
 \widetilde X\times_{\mathbf A^1}\widetilde X;
\]
\item a coefficient carrier
\[
 \widetilde K\in\Dctf(\widetilde X,A)
\]
and a cohomological correspondence
\[
 \widetilde u:L\widetilde c_2^*\widetilde K
 \longrightarrow R\widetilde c_1^!\widetilde K
 \quad\text{in }D(\widetilde C,A);
\]
\item over \(\mathbf G_m\), identifications
\[
 \widetilde C_{\mathbf G_m}\simeq C\times\mathbf G_m,
 \quad d_{\mathbf G_m}\simeq c\times\operatorname{id}_{\mathbf G_m},
 \quad \widetilde K_{\mathbf G_m}\simeq
 \operatorname{pr}_X^*K,
 \quad \widetilde u_{\mathbf G_m}\simeq
 u\times\operatorname{id}_{\mathbf G_m},
\]
compatible with both legs and with exceptional exchange.
\end{enumerate}

For \(t\in\{0,1\}\), put
\[
 X_t=X\times\{t\},\qquad
 C_t:=\widetilde C\times_{\mathbf A^1}\{t\}.
\]
Let \(\iota_t:X_t\hookrightarrow\widetilde X\) and
\(\jmath_t:C_t\hookrightarrow\widetilde C\) be the fiber immersions.  The
fiber data are
\[
 c_{t,r}:C_t\longrightarrow X_t,
 \qquad
 K_t=L\iota_t^*\widetilde K,
\]
and
\[
\begin{aligned}
 u_t:\quad Lc_{t,2}^*K_t
 &\xrightarrow{\sim}L\jmath_t^*L\widetilde c_2^*\widetilde K
 \xrightarrow{L\jmath_t^*\widetilde u}
 L\jmath_t^*R\widetilde c_1^!\widetilde K\\
 &\xrightarrow{\sim}Rc_{t,1}^!K_t.
\end{aligned}
\]
The fiber at \(1\) is identified with \((C,c,K,u)\).  At \(0\), the map
\(d_0:C_0\to\mathfrak D_0=C_{\Delta_X}(X\times X)\), the two induced
legs \(c_{0,1},c_{0,2}:C_0\to X\), the carrier \(K_0\), and the arrow
\(u_0\) form the normal-cone correspondence.
\iffalse
\[
 c_0=(c_{0,1},c_{0,2}),\qquad K_0,\qquad
 u_0:Lc_{0,2}^*K_0\longrightarrow Rc_{0,1}^!K_0.
\]

\fi
Let \(\Delta_{\widetilde X/\mathbf A^1}\) be the relative diagonal and
define the fixed family by
\[
 \widetilde F:=
 \widetilde C\times_{\widetilde X\times_{\mathbf A^1}\widetilde X,\,
 \widetilde c,\,\Delta_{\widetilde X/\mathbf A^1}}\widetilde X.
\]
Write \(i_{\widetilde F}:\widetilde F\to\widetilde C\) and
\(d_{\widetilde F}:\widetilde F\to\widetilde X\) for its projections.
For \(t\in\{0,1\}\), define
\[
 F_t=\widetilde F\times_{\mathbf A^1}\{t\}.
\]
Base change gives canonical identifications
\[
 F_t\xrightarrow{\sim}\Fix(c_t),
 \qquad c_{t,1}i_{F_t}=d_{F_t}=c_{t,2}i_{F_t}.
\]
When
\(\widetilde F\to\mathbf A^1\) is proper, the trivialization over
\(\mathbf G_m\), followed by nearby-cycle specialization at zero, defines
\[
 \operatorname{sp}_{\mathrm{DNC}}:
 \Hom_{D(F_1,A)}(A_{F_1},\omega_{F_1,A})
 \longrightarrow
 \Hom_{D(F_0,A)}(A_{F_0},\omega_{F_0,A}).
\]
Under the displayed fiber identifications, \(F_1=\Fix(c)\) and
\(F_0=\Fix(c_0)\).
\end{definition}

\begin{proposition}[Specialization of the trace class]
For a deformation datum with proper fixed family,
\[
 \operatorname{sp}_{\mathrm{DNC}}
 \bigl(\operatorname{Tr}_c(u)\bigr)=
 \operatorname{Tr}_{c_0}(u_0).
\]
The proper-family trace identifies their integrals:
\[
 \int_{\Fix(c)}\operatorname{Tr}_c(u)=
 \int_{\Fix(c_0)}\operatorname{Tr}_{c_0}(u_0).
\]
This equality is compatible with distinguished triangles and proper
pushforward.
\depends{Nearby-cycle specialization for the deformation space, proper
base change for the stated proper fixed family, purity for the diagonal,
and functoriality of trace classes.}
\end{proposition}

\begin{proposition}[Simple-point deformation package]
Let \(h:X\to X\) fix \(z\), assume that \(X\) is smooth and \(K\) is
lisse in an \(h\)-stable neighborhood of \(z\), and assume that
\[
 1-dh_z:T_zX\longrightarrow T_zX
\]
is invertible.  After etale excision around \(z\), the graph--diagonal
trace problem admits a deformation datum whose fixed family is the section
\(z\times\mathbf A^1\), hence is proper over \(\mathbf A^1\).  Its special
fiber is the linear graph correspondence of
\(dh_z:T_zX\to T_zX\); its fixed scheme is the zero section because
\(1-dh_z\) is invertible.  The specialized coefficient object is the
constant perfect stalk \(K_z\) on the normal fiber.
\depends{Etale excision, transversality of the graph and diagonal, the
normal-cone construction for a transverse intersection, and lisse
trivialization of the coefficient complex near the fixed point.}
\end{proposition}

\begin{proposition}[Vector-bundle trace identity]
Let \(p:V\to Z\) be a vector bundle of rank \(d\), with zero section
\(z:Z\to V\).  Purity identifies
\[
 Rz^!A_V\simeq A_Z(-d)[-2d],
\qquad
 Rp_!A_V\simeq A_Z(-d)[-2d].
\]
Under these identifications, the structural trace
\[
 Rp_!Rp^!A_Z\longrightarrow A_Z
\]
is normalized by the Thom class.  The trace class of the normal-bundle
correspondence is identified with the trace class of the identity
correspondence on \(Z\), and its geometric fibers are the identity
correspondence on a point.
\depends{Smooth purity, homotopy invariance, the Thom isomorphism, and
normalization of the exceptional counit.}
\end{proposition}

\begin{lemma}[Identity correspondence on a point]
Let \(P\in\Perf(A)\) and \(v:P\to P\).  For the identity correspondence of
\(\Spec\Omega\), the Verdier local term equals the categorical trace:
\[
 \operatorname{LT}_{\Spec\Omega}(v)=\operatorname{tr}_A(v\mid P).
\]
\depends{The definitions of evaluation, coevaluation, trace class, and
integration for the identity map.}
\end{lemma}

\begin{proposition}[The six-functor trace diagram]
Let \(P=\RGammac(X,K)\), and assume that \(c_2\) and
\(F=\Fix(c)\) are proper.  The global row
\[
 A\xrightarrow{\operatorname{coev}_P}P\otimes_A^LP^\vee
 \xrightarrow{T_c(u)\otimes1}P\otimes_A^LP^\vee
 \xrightarrow{\mathrm{sym}}P^\vee\otimes_A^LP
 \xrightarrow{\operatorname{ev}_P}A
\]
and the fixed-point row
\[
 A\longrightarrow R\Gamma(F,A_F)
 \xrightarrow{R\Gamma(F,\operatorname{Tr}_c(u))}
 R\Gamma(F,\omega_{F,A})
 \xrightarrow{\varepsilon_{a_F}}A
\]
have the same composite in \(\End_{D(A)}(A)\).  The first composite is
\(\operatorname{tr}_A(T_c(u)\mid P)\), and the second is
\(\int_F\operatorname{Tr}_c(u)\).
\depends{Compact-support Kunneth, Verdier biduality, the projection formula,
proper base change for \(c_2\), exceptional base change along the diagonal,
and the triangle identities for the ordinary and exceptional adjunctions.}
\end{proposition}

\begin{theorem}[Global Lefschetz--Verdier trace comparison]
Let \(c=(c_1,c_2):C\to X\times X\), let
\(K\in\Dctf(X,A)\), and let
\[
 u:Lc_2^*K\longrightarrow Rc_1^!K.
\]
Assume \(c_2\) and \(\Fix(c)\to s_\Omega\) are proper.  Then
\[
 \operatorname{tr}_A
 \bigl(T_c(u)\mid\RGammac(X,K)\bigr)
 =
 \int_{\Fix(c)}\operatorname{Tr}_c(u)
 \quad\text{in }A.
\]
\depends{Perfectness of compact-support cohomology and the six-functor trace
diagram.}
\end{theorem}

\begin{corollary}[Lefschetz--Verdier component formula]
Under the hypotheses of the global comparison,
\[
 \operatorname{tr}_A
 \bigl(T_c(u)\mid\RGammac(X,K)\bigr)
 =
 \sum_{\beta\in\pi_0(\Fix(c))}
 \operatorname{LT}_\beta(u).
\]
\depends{The global comparison and additivity of integration on the finite
set of open-and-closed connected components.}
\end{corollary}

\section{The Frobenius graph and its local calculation}

\begin{definition}[Graph correspondence]
For a proper endomorphism \(h:X\to X\) and a morphism
\(v:Lh^*K\to K\), define
\[
 C=X,\qquad c_1=\operatorname{id}_X,\qquad c_2=h,\qquad u=v.
\]
Then
\[
 \Fix(c)=\Fix(h)
\]
is the scheme-theoretic equalizer of \(h\) and \(\operatorname{id}_X\).
\end{definition}

\begin{proposition}[Graph construction of the global operator]
For a graph correspondence \((h,v)\), the correspondence operator
\(T_c(v)\) is the composite
\[
 \RGammac(X,K)\longrightarrow\RGammac(X,Lh^*K)
 \xrightarrow{\RGammac(X,v)}\RGammac(X,K),
\]
where the first arrow is induced by the unit
\(K\to Rh_*Lh^*K\) and the identification \(Rh_*\simeq Rh_!\).
It is functorial in \((K,v)\), additive in distinguished triangles, and
compatible with composition of proper endomorphisms.
\depends{The definition of the global correspondence operator and
properness of \(h\).}
\end{proposition}

For \(m\geq1\), apply this definition to
\[
 h=\Phi^m,\qquad
 u_m=\varphi^{[m]}\circ\kappa_K^{(m)}:
 L(\Phi^m)^*K\longrightarrow K.
\]

\begin{proposition}[Comparison of the two global Frobenius operators]
The graph operator \(T_{\Phi^m}(u_m)\) is equal to
\[
 F_{c,K}^{\,m}:\RGammac(X,K)\to\RGammac(X,K)
\]
defined from \(g_\gamma^m\) and \(\varphi^{[m]}\).  These equalities are
natural in strict Weil complexes and compatible with the cocycle law in
\(m\).
\depends{The comparison \(\kappa^{(m)}\), compact-support base change for
the automorphism \(g_\gamma^m\), and the graph construction.}
\end{proposition}

\begin{definition}[Simple fixed point]
A geometric fixed point \(z\) of \(h:X\to X\) is simple if
\(\Fix(h)\to\Spec\Omega\) is etale at \(z\).  When the Zariski tangent
space is used and \(X\) is smooth at \(z\), this is equivalent to
invertibility of
\[
 1-dh_z:T_zX\longrightarrow T_zX.
\]
At a simple fixed point, the local fixed scheme is \(\Spec\Omega\) and has
the multiplicity-one orientation supplied by etale purity.
\end{definition}

\begin{theorem}[Scheme-theoretic Frobenius fixed points]
For every \(m\geq1\),
\[
 \Fix(\Phi^m)\xrightarrow{\sim}
 \coprod_{y\in X_0(\mathbf F_{q^m})}\Spec\bar k.
\]
In particular, \(\Fix(\Phi^m)\) is finite etale over \(\bar k\), and every
fixed point is simple, without a smoothness assumption.  On a smooth stratum,
one has \(d\Phi^m=0\), so \(1-d\Phi^m\) is the identity.  The graph and
diagonal therefore meet transversally there with local multiplicity one.  The
same local normalization holds after stable stratification and excision on
\(X\).
\depends{The affine description of scheme Frobenius, finiteness of the
fixed equalizer, and finite-field descent for rational points.}
\end{theorem}

\begin{proposition}[Smooth lisse Frobenius-stable stratification]
The universal homeomorphism \(X_{\mathrm{red}}\to X\) preserves the
Frobenius trace class and all local terms.  The reduced scheme has a finite
stratification by smooth locally closed subschemes defined over \(\Fq\),
refined so that every cohomology sheaf of \(K\) is lisse on each stratum;
each stratum is stable under \(\Phi^m\).  The Frobenius trace class and its
local terms are the additive extensions, under stable localization and
excision, of the corresponding trace classes on these smooth lisse strata.
\depends{Topological invariance, existence of a smooth stratification over
the perfect field \(\Fq\), constructible devissage, extension by zero,
stable localization, and etale excision.}
\end{proposition}

\begin{proposition}[Local trace at a simple graph fixed point]
Let \(h:X\to X\) be proper, let \(z\) be a simple fixed point, and let
\(v:Lh^*K\to K\).  Assume that \(X\) is smooth and \(K\) is lisse in an
\(h\)-stable neighborhood of \(z\).  The restriction of the Verdier trace
class to \(z\) is the categorical trace of the induced stalk endomorphism
\[
 v_z:K_z\xrightarrow{\sim}(Lh^*K)_z\longrightarrow K_z.
\]
Consequently,
\[
 \operatorname{LT}_z(v)=\operatorname{tr}_A(v_z\mid K_z).
\]
\depends{The simple-point deformation package, specialization of the trace
class, the Thom reduction, and the identity-correspondence calculation on a
point.}
\end{proposition}

\begin{proposition}[Frobenius local-term formula]
For \(y\in X_0(\mathbf F_{q^m})\), viewed as a component of
\(\Fix(\Phi^m)\),
\[
 \operatorname{LT}_y(u_m)=
 \operatorname{tr}_A(F_{y,m}\mid K_{\bar y}).
\]
If \(y\) lies over a closed point \(x\) of degree \(d_x\mid m\), this is
\[
 \operatorname{tr}_A
 \bigl(F_x^{\,m/d_x}\mid K_{\bar x}\bigr).
\]
\depends{The scheme-theoretic fixed-point theorem, reduction to smooth lisse
Frobenius-stable strata, the simple-point local trace theorem, extension by
zero, and the scheme/Galois Frobenius comparison on stalks.}
\end{proposition}

\begin{theorem}[Grothendieck--Lefschetz formula over a finite ring]
Let \((K,\varphi)\in\DctfW(X,A)\).  Assume that the structural maps and
fixed-point base changes used in the trace-class construction are
\(A\)-exceptionally admissible with the required composition and exchange
data.  For every \(m\geq1\),
\[
 \boxed{
 \operatorname{tr}_A
 \bigl(F_{c,K}^{\,m}\mid\RGammac(X,K)\bigr)
 =
 \sum_{y\in X_0(\mathbf F_{q^m})}
 \operatorname{tr}_A(F_{y,m}\mid K_{\bar y}).}
\]
Equivalently,
\[
 \operatorname{tr}_A
 \bigl(F_{c,K}^{\,m}\mid\RGammac(X,K)\bigr)
 =
 \sum_{\substack{x\in|X_0|\\d_x\mid m}}
 d_x\,
 \operatorname{tr}_A
 \bigl(F_x^{\,m/d_x}\mid K_{\bar x}\bigr).
\]
Both sides lie in \(A\), and both sums are finite.
\depends{The global Lefschetz--Verdier comparison, graph construction,
global Frobenius comparison, the fixed-point theorem, and the Frobenius
local-term formula.}
\end{theorem}

\section{Finite coefficient change}

\begin{theorem}[Compatibility with a finite coefficient map]
Let \(\rho:A\to B\) be a homomorphism of finite commutative
\(\Zl\)-algebras.  Derived scalar extension
\[
 -\otimes_A^LB:D(Y,A)\to D(Y,B)
\]
preserves \(\Dctf\).  For every separated finite-type \(f:Y\to Z\), there
are natural isomorphisms
\[
\begin{aligned}
(Rf_!K)\otimes_A^LB&\xrightarrow{\sim}Rf_!(K\otimes_A^LB),\\
(Lf^*L)\otimes_A^LB&\xrightarrow{\sim}Lf^*(L\otimes_A^LB).
\end{aligned}
\]
If \(f\) is exceptionally admissible for both coefficient rings and the
two exceptional adjunction data are compatible with scalar extension, there
is in addition a natural isomorphism
\[
 (Rf_A^!L)\otimes_A^LB
 \xrightarrow{\sim}Rf_B^!(L\otimes_A^LB).
\]
These maps commute with the available units, counits, composition,
compact-support compactifications, and a composite \(A\to B\to C\).

For a Weil complex \(K_A\), they give Frobenius-equivariant isomorphisms
\[
\begin{aligned}
\RGammac(X,K_A)\otimes_A^LB&\xrightarrow{\sim}\RGammac(X,K_B),\\
(K_A)_{\bar y}\otimes_A^LB&\xrightarrow{\sim}(K_B)_{\bar y}.
\end{aligned}
\]
For a cohomological correspondence \(u_A\), they identify
\[
 \operatorname{Tr}_c(u_A)\otimes_A^LB
 =\operatorname{Tr}_c(u_B),
\qquad
 \rho\bigl(\operatorname{LT}_\beta(u_A)\bigr)
 =\operatorname{LT}_\beta(u_B).
\]
Consequently the finite-ring trace formula over \(A\) maps to the
finite-ring trace formula over \(B\).
\depends{Coefficient change for lower shriek and pullback; for the
exceptional comparison, admissibility over both rings and compatibility of
the chosen adjoints; Verdier duality, categorical trace, and the fixed-point
evaluation morphism.}
\end{theorem}

\section{Adic towers and reconstruction}

Let \(E/\Ql\) be finite, let \(\OO_E\) be its ring of integers, choose a
uniformizer \(\varpi\), and put
\[
 A_n:=\OO_E/\varpi^n\OO_E.
\]

\begin{definition}[Normalized adic coefficient system]
An \(A_\bullet\)-complex on \(Y\) is a family
\[
 K_\bullet=(K_n,t_{n+1,n})_{n\geq1},\qquad
 K_n\in D(Y,A_n),
\]
with transition morphisms compatible with the quotient maps.  It is
normalized if, for \(m\leq n\),
\[
 K_n\otimes_{A_n}^LA_m\xrightarrow{\sim}K_m.
\]
Let
\[
 \mathcal D^b_{\mathrm{ctf,adic}}(Y,\OO_E)
\]
be the full triangulated subcategory of the derived category of systems
consisting of normalized systems whose levels lie in \(\Dctf(Y,A_n)\),
with one cohomological bound and one Tor-amplitude valid for all \(n\).

A strict adic Weil system is such a system represented by a cochain complex
of systems together with a cochain isomorphism from its
\(g_\gamma\)-pullback to itself.  Its triangulated category is denoted
\[
 \mathcal D^b_{\mathrm{ctf,W,adic}}(X,\OO_E).
\]
\end{definition}

\begin{definition}[Derived adic realization]
Let
\[
 \underline{\OO_E}_Y:=\varprojlim_n(A_n)_Y.
\]
For every \(n\), let
\[
 \operatorname{Res}_{n,Y}:D(Y,A_n)\longrightarrow
 D(Y,\underline{\OO_E}_Y)
\]
be restriction of scalars along
\(\underline{\OO_E}_Y\twoheadrightarrow(A_n)_Y\).  All inverse limits of
level objects below are formed after applying this functor.
For a normalized system \(K_\bullet\), define
\[
 K_{\OO_E}:=R\!\varprojlim_n\operatorname{Res}_{n,Y}K_n
 \quad\text{in }D(Y,\underline{\OO_E}_Y).
\]
An \(\underline{\OO_E}_Y\)-complex \(L\) is derived
\(\varpi\)-complete if
\[
 L\xrightarrow{\sim}
 R\!\varprojlim_n(L\otimes_{\OO_E}^LA_n).
\]
\end{definition}

\begin{proposition}[Levelwise compact support and stalk compatibility]
For
\[
 K_\bullet\in
 \mathcal D^b_{\mathrm{ctf,W,adic}}(X,\OO_E),
\]
the systems
\[
 P_n:=\RGammac(X,K_n),\qquad
 P_{\bar y,n}:=(K_n)_{\bar y}
\]
are normalized systems of perfect \(A_n\)-complexes.  Every global and local
Frobenius operator is compatible with the transition maps.  Define their
realizations in \(D(\OO_E)\) by
\[
 P_{\OO_E}:=
 R\!\varprojlim_n\operatorname{Res}_{n,s}P_n,\qquad
 P_{\bar y,\OO_E}:=
 R\!\varprojlim_n\operatorname{Res}_{n,s}P_{\bar y,n}.
\]
There are Frobenius-equivariant realization isomorphisms
\[
\begin{aligned}
 \chi_c(K_\bullet):\quad
 \RGammac(X,K_{\OO_E})
 &\xrightarrow{\sim}P_{\OO_E},\\
 \chi_{\bar y}(K_\bullet):\quad
 (K_{\OO_E})_{\bar y}
 &\xrightarrow{\sim}P_{\bar y,\OO_E}.
\end{aligned}
\]
Here
\[
 \RGammac(X,K_{\OO_E})
 :=R\Gamma\bigl(s,Ra_!K_{\OO_E}\bigr),
 \qquad
 (K_{\OO_E})_{\bar y}:=L\bar y^*K_{\OO_E}.
\]
Both realized objects belong to \(\Perf(\OO_E)\), and
\[
\begin{aligned}
P_{\OO_E}\otimes_{\OO_E}^LA_n&\xrightarrow{\sim}P_n,\\
P_{\bar y,\OO_E}\otimes_{\OO_E}^LA_n&\xrightarrow{\sim}P_{\bar y,n}
\end{aligned}
\]
are Frobenius-equivariant.
\depends{Finite-level coefficient change, uniform boundedness, derived
completeness, and finiteness of compactly supported direct image.}
\end{proposition}

\begin{proposition}[Trace compatibility in the adic tower]
Let \(F_{\OO_E}\) be a global or local Frobenius endomorphism on one of the
realized perfect complexes above, and let \(F_n\) be its reduction.  Then
\[
 \operatorname{tr}_{\OO_E}(F_{\OO_E})
 \bmod\varpi^n
 =
 \operatorname{tr}_{A_n}(F_n).
\]
The resulting elements of \(A_n\) form a compatible inverse system.
\depends{Perfectness of the realized complexes and scalar compatibility of
categorical trace.}
\end{proposition}

\begin{proposition}[Adic reconstruction of a trace identity]
Let \(\alpha,\beta\in\OO_E\).  If their images in \(A_n\) are equal for
every \(n\geq1\), then \(\alpha=\beta\).  Hence a compatible family of
finite-level trace formulas reconstructs a unique equality in \(\OO_E\).
Explicitly, the identification
\[
 \OO_E\xrightarrow{\sim}\varprojlim_nA_n
\]
sends the integral global trace and each integral local trace to their
finite-level trace systems.
\depends{\(\varpi\)-adic completeness and separatedness of \(\OO_E\), and
trace compatibility in the adic tower.}
\end{proposition}

\begin{theorem}[Integral adic trace formula for a Frobenius correspondence]
Let
\[
 L_\bullet\in
 \mathcal D^b_{\mathrm{ctf,adic}}(X,\OO_E)
\]
and fix \(m\geq1\).  A compatible integral Frobenius correspondence is a
morphism of normalized systems
\[
 v_\bullet:L(\Phi^m)^*L_\bullet\longrightarrow L_\bullet.
\]
Assume that the finite-level exceptional adjoints and Verdier-duality data
used by the graph trace formula are compatible with reduction and satisfy
the uniform amplitude and Tor bounds required for adic reconstruction.
It induces an endomorphism
\(T_{\Phi^m}(v_{\OO_E})\) of
\(P_{\OO_E}=\RGammac(X,L_{\OO_E})\), and, for every
\(y\in X_0(\mathbf F_{q^m})\), an endomorphism
\(v_{y,\OO_E}\) of the perfect stalk
\(P_{\bar y,\OO_E}\).  Then
\[
 \operatorname{tr}_{\OO_E}
 \bigl(T_{\Phi^m}(v_{\OO_E})\mid P_{\OO_E}\bigr)
 =
 \sum_{y\in X_0(\mathbf F_{q^m})}
 \operatorname{tr}_{\OO_E}
 \bigl(v_{y,\OO_E}\mid P_{\bar y,\OO_E}\bigr).
\]
No invertibility of \(v_\bullet\) is required.  Reduction modulo
\(\varpi^n\) is the finite-ring Lefschetz--Verdier formula for the graph
correspondence \(v_n:L(\Phi^m)^*L_n\to L_n\).
\depends{The finite-ring Lefschetz--Verdier formula for an arbitrary
cohomological correspondence at every level, compatibility with reduction,
and adic reconstruction of traces.}
\end{theorem}

\begin{corollary}[Integral adic Grothendieck--Lefschetz formula]
For
\[
 K_\bullet\in
 \mathcal D^b_{\mathrm{ctf,W,adic}}(X,\OO_E)
\]
and \(m\geq1\),
\[
 \operatorname{tr}_{\OO_E}
 \bigl(F_{c,K_{\OO_E}}^{\,m}\mid P_{\OO_E}\bigr)
 =
 \sum_{y\in X_0(\mathbf F_{q^m})}
 \operatorname{tr}_{\OO_E}
 \bigl(F_{y,m,K_{\OO_E}}\mid P_{\bar y,\OO_E}\bigr).
\]
Reduction modulo \(\varpi^n\) is exactly the finite-ring formula for
\(K_n\).
\depends{The integral adic correspondence formula applied to the iterated
strict Weil arrow and levelwise Frobenius compatibility.}
\end{corollary}

\section{Finite extensions of \texorpdfstring{\(\Ql\)}{Qell}}

\begin{definition}[Rational constructible Weil category]
Put
\[
 \underline E_Y:=\underline{\OO_E}_Y\otimes_{\OO_E}E.
\]
Let \(D(Y,E)\) be the derived category of sheaves of
\(\underline E_Y\)-modules on \(Y_{\mathrm{pro\acute et}}\), and let
\(\Dbc(Y,E)\) be its full triangulated subcategory of bounded complexes with
constructible, finite-dimensional \(E\)-stalks.

Let \(\mathsf{Sh}_{\mathrm W}(X,E)\) be the abelian category of sheaves
with a strict isomorphism \(g_\gamma^*\mathcal F\simeq\mathcal F\).
Define
\[
 \DbcW(X,E)\subset
 D\bigl(\mathsf{Sh}_{\mathrm W}(X,E)\bigr)
\]
by requiring the underlying complex to lie in \(\Dbc(X,E)\).
\end{definition}

\begin{definition}[Rational exceptional-adjoint admissibility]
A separated finite-type morphism \(f:Y\to Z\) is
\emph{\(E\)-exceptionally admissible} if, on compatible finite and adic
coefficient presentations, the functors \(Rf_!\) have right adjoints with
reduction-compatible units and counits, a uniform finite cohomological
amplitude, and uniform finite Tor-amplitude on constructible objects.  The
derived-complete reconstruction of this system is required to preserve
these adjunction and composition identities.  Only under this hypothesis is
the reconstructed rational right adjoint denoted
\[
 Rf^!:D(Z,E)\longrightarrow D(Y,E).
\]
\end{definition}

\begin{theorem}[Lattice comparison and rational perfectness]
Sheafwise realization and inversion of \(\varpi\) induce an exact
equivalence
\[
 \Kar\!\left(
 \mathcal D^b_{\mathrm{ctf,adic}}(Y,\OO_E)[1/\varpi]\right)
 \xrightarrow{\sim}\Dbc(Y,E).
\]
For every separated finite-type \(f:Y\to Z\), lower shriek gives an exact
functor
\[
 Rf_!:\Dbc(Y,E)\longrightarrow\Dbc(Z,E).
\]
If \(f\) is \(E\)-exceptionally admissible, there is in addition an exact
right adjoint
\[
 Rf^!:\Dbc(Z,E)\longrightarrow\Dbc(Y,E),
 \qquad Rf_!\dashv Rf^!.
\]
Compact support and geometric stalks take values in
\(\Perf(E)\).  Explicitly, every \(K_E\in\Dbc(Y,E)\) is represented by
\[
 L_\bullet\in\mathcal D^b_{\mathrm{ctf,adic}}(Y,\OO_E),
 \qquad e^2=e\in
 \End\bigl(L_\bullet[1/\varpi]\bigr),
\]
and a splitting after scalar extension
\[
 K_E\xrightarrow{\iota}L_E
 \xrightarrow{p}K_E,\qquad p\iota=1,\quad\iota p=e,
 \qquad L_E:=L_\bullet\otimes_{\OO_E}E.
\]
Compact support and stalks of \(K_E\) are the corresponding idempotent
summands of those of \(L_E\).  This presentation concerns the underlying
geometric complex only; it does not assert that a chosen strict Weil
isomorphism preserves an \(\OO_E\)-lattice.
\depends{Adic realization, derived completeness, and idempotent completion
after inverting \(\varpi\); for \(Rf^!\), the finite/adic admissibility
datum, its uniform amplitude and Tor bounds, and reconstruction of the
adjunction.}
\end{theorem}

\begin{proposition}[Clearing denominators in a strict Weil arrow]
Let \((K_E,\varphi_E)\in\DbcW(X,E)\), choose a presentation
\((L_\bullet,e,\iota,p)\) as above, and use the pulled-back splitting on
\(Lg_\gamma^*L_E\).  There exist a nonzero \(a\in\OO_E\) and an integral
morphism
\[
 w:Lg_\gamma^*L_\bullet\longrightarrow L_\bullet
\]
whose scalar extension is the following block morphism, equal to
\(a\varphi_E\) on the summand \(K_E\) and zero on its complement:
\[
 w\otimes E=
 \iota\circ(a\varphi_E)\circ Lg_\gamma^*p.
\]
The morphism \(w\) need not be invertible.  The inverse
\(\varphi_E^{-1}\) can be cleared with a separate denominator, but is not
needed below.  Define
\[
 w^{[m]}:=w\circ Lg_\gamma^*w\circ\cdots\circ
 L(g_\gamma^{m-1})^*w.
\]
On the summand \(K_E\), its scalar extension is
\[
 p\circ(w^{[m]}\otimes E)\circ L(g_\gamma^m)^*\iota
 =a^m\varphi_E^{[m]}.
\]
If \(v_m=w^{[m]}\circ\kappa_{L}^{(m)}\) and
\(u_m=\varphi_E^{[m]}\circ\kappa_{K}^{(m)}\), then the restriction of
\(v_m\otimes E\) to \(K_E\) is \(a^mu_m\), while its complementary
summand is zero.
\depends{The description of morphisms in the localization obtained by
inverting \(\varpi\), the Karoubi splitting, pullback compatibility with
scalars and idempotents, and the strict Weil iterate cocycle.}
\end{proposition}

\begin{theorem}[Grothendieck--Lefschetz formula over \(E\)]
Assume that the structural and fixed-point morphisms used in Verdier
duality are \(E\)-exceptionally admissible.  For
\((K_E,\varphi_E)\in\DbcW(X,E)\) and \(m\geq1\),
\[
 \operatorname{tr}_E
 \bigl(F_{c,K_E}^{\,m}\mid\RGammac(X,K_E)\bigr)
 =
 \sum_{y\in X_0(\mathbf F_{q^m})}
 \operatorname{tr}_E
 \bigl(F_{y,m,K_E}\mid(K_E)_{\bar y}\bigr).
\]
Equivalently,
\[
 \sum_i(-1)^i\Tr
 \bigl(F_{c,K_E}^{\,m}\mid H_c^i(X,K_E)\bigr)
 =
 \sum_{y\in X_0(\mathbf F_{q^m})}
\sum_j(-1)^j\Tr
\bigl(F_{y,m,K_E}\mid H^j((K_E)_{\bar y})\bigr).
\]
For every denominator-cleared lift \((a,w)\), scalar extension of its
integral correspondence identity is \(a^m\) times the displayed equality.
No Frobenius-stable integral lattice is part of the hypotheses.
\depends{A Karoubi presentation of the underlying complex, denominator
clearing for the Weil arrow and its iterates, the integral adic trace
formula for a noninvertible Frobenius correspondence, scalar extension
\(\OO_E\to E\), and additivity and \(E\)-linearity of categorical trace.}
\end{theorem}

\begin{proposition}[Finite-extension coefficient coherence]
For an inclusion \(E\hookrightarrow E'\) of finite extensions of \(\Ql\),
derived scalar extension gives an exact triangulated functor
\[
 -\otimes_EE':\DbcW(X,E)\to\DbcW(X,E').
\]
It commutes with compact support, stalks, trace classes, local terms, and
every Frobenius iterate.  The identity and composition isomorphisms
\[
 K_E\otimes_EE\simeq K_E,\qquad
 (K_E\otimes_EE')\otimes_{E'}E''\simeq K_E\otimes_EE''
\]
satisfy the unit and associativity coherence conditions.  Traces map under
\(E\hookrightarrow E'\) to the corresponding \(E'\)-valued traces.
Compatible exceptional-adjoint data with the same uniform bounds base
change from \(E\) to \(E'\), so exceptional admissibility is preserved
whenever these comparison data are supplied.
\depends{Finite coefficient change, lattice comparison, scalar
compatibility of categorical trace, and base change of the supplied
exceptional-adjoint data.}
\end{proposition}

\section{Algebraically closed coefficients}

Fix \(\Qlb\), and let \(\mathcal I_\ell\) be the filtered category of finite
extensions \(E/\Ql\) contained in \(\Qlb\).

\begin{definition}[The \(\Qlb\)-coefficient Weil category]
Form the filtered triangulated 2-colimit
\[
 \mathcal C_{\Qlb}^{0}:=
 2\mathop{\mathrm{colim}}_{E\in\mathcal I_\ell}\DbcW(X,E).
\]
An object is represented by \((E,K_E)\).  A morphism is represented, after
passing to a common finite overfield \(M\), by a morphism
\[
 K_E\otimes_EM\longrightarrow L_{E'}\otimes_{E'}M,
\]
with two representatives identified after a further finite extension.
Distinguished triangles are those represented at a finite stage.  Define
\[
 \DbcW(X,\Qlb):=\Kar(\mathcal C_{\Qlb}^{0}).
\]
Thus a general object is represented, after enlarging \(E\), by
\[
 (E,K_E,e),\qquad e^2=e,
\]
where \(e\) is Frobenius-equivariant.
\end{definition}

\begin{definition}[Trace on a Karoubi summand]
Let \(P\in\Perf(E)\), let \(e:P\to P\) be idempotent, and let
\(F:P\to P\) commute with \(e\).  After scalar extension to \(\Qlb\), choose
a splitting
\[
 Q\xrightarrow{i}P_{\Qlb}\xrightarrow{p}Q,\qquad
 p i=1_Q,\qquad i p=e_{\Qlb}.
\]
The induced endomorphism is
\[
 F_Q:=pF_{\Qlb}i.
\]
Define
\[
 \operatorname{tr}_{\Qlb}(F\mid(P,e))
 :=
 \operatorname{tr}_{\Qlb}(F_Q\mid Q)
 =
 \operatorname{tr}_{\Qlb}(F_{\Qlb}e_{\Qlb}\mid P_{\Qlb}).
\]
The value is independent of the splitting.
\end{definition}

\begin{proposition}[Karoubi trace and representative independence]
Let \((E,K_E,e)\) represent an object of \(\DbcW(X,\Qlb)\).  Compact support
and every geometric stalk carry \(e\) to Frobenius-equivariant idempotents.
Their images are perfect \(\Qlb\)-complexes, and the preceding definition
gives their Frobenius traces.  If another triple represents the same object,
then after passage to a common finite overfield the two idempotents are
conjugate.  The compact-support traces and all local traces therefore agree.
These comparisons are transitive and commute with every Frobenius iterate.
\depends{Finite-extension coherence, idempotent completeness of perfect
complexes over a field, conjugacy invariance, and additivity of trace.}
\end{proposition}

\begin{proposition}[Trace formula on an equivariant summand]
Let \((E,K_E,e)\) be as above, and assume the exceptional-adjoint
admissibility hypotheses of the \(E\)-valued trace formula.  For every
\(m\geq1\),
\[
\begin{aligned}
&\operatorname{tr}_E
\bigl(F_{c,K_E}^{\,m}e_c\mid\RGammac(X,K_E)\bigr)\\
&\qquad =
\sum_{y\in X_0(\mathbf F_{q^m})}
\operatorname{tr}_E
\bigl(F_{y,m,K_E}e_{\bar y}\mid(K_E)_{\bar y}\bigr).
\end{aligned}
\]
After scalar extension, this is the trace formula on the corresponding
\(\Qlb\)-valued Karoubi summands.
\depends{The \(E\)-valued trace formula applied to the endomorphism of the
distinguished decomposition defined by \(e\), and additivity of trace.}
\end{proposition}

\begin{theorem}[Grothendieck--Lefschetz formula over \(\Qlb\)]
Let \((K,\varphi)\in\DbcW(X,\Qlb)\).  Assume that at one finite coefficient
field of definition the structural and fixed-point morphisms have compatible
exceptional-adjoint data; by finite-extension coherence this hypothesis is
independent of enlarging that field.  For every \(m\geq1\),
\[
 \boxed{
 \sum_i(-1)^i\Tr
 \bigl(F_{c,K}^{\,m}\mid H_c^i(X,K)\bigr)
 =
 \sum_{y\in X_0(\mathbf F_{q^m})}
 \sum_j(-1)^j\Tr
 \bigl(F_{y,m,K}\mid H^j(K_{\bar y})\bigr).}
\]
All vector spaces are finite-dimensional and only finitely many degrees
occur.  Every trace descends to a finite coefficient field, and the equality
is independent of the finite field, the representative, the common
overfield, and the splitting of the idempotent.
\depends{Finite-stage descent, the trace formula on equivariant summands,
and representative independence.}
\end{theorem}

\begin{corollary}[Point-counting trace formula]
For the constant object \(\underline{\Qlb}_X\) in degree zero with its canonical strict
Weil structure, every local Frobenius is the identity on a one-dimensional
stalk, and
\[
 \#X_0(\mathbf F_{q^m})
 =
 \sum_i(-1)^i\Tr
 \bigl(F_c^m\mid H_c^i(X,\Qlb)\bigr).
\]
The same formula holds over every finite extension \(E/\Ql\), and its
images under coefficient extension agree.
\depends{The \(\Qlb\)-valued trace formula and the normalization of the
constant Weil sheaf at a rational point.}
\end{corollary}

\section{References}

The finite-coefficient trace-class formalism, its compatibility with proper
pushforward and localization, and deformation-to-normal-cone compatibility are the content
of the Lefschetz--Verdier formalism in SGA 5 and the trace report in SGA
4 1/2.  The strict Weil convention and the passage to rational coefficients
agree with the conventions of \emph{Weil I} and \emph{Weil II}.  The
pro-etale coefficient categories and topological invariance may be compared
with the corresponding parts of the Stacks Project and the pro-etale
formalism of Bhatt--Scholze.

\end{document}
